% Source: https://artofproblemsolving.com/community/c3h1429460_geo_bee_8
% Stable identity: aops-topic-1429460-q08
Condition on the set of six chosen points; label them $1,\dots,6$ in circular order. Every way of choosing which $3$ of them are $A,B,C$ is equally likely, so there are $\binom{6}{3}=20$ equally likely splits into $\{A,B,C\}$ and $\{D,E,F\}$.

\begin{itemize}
\item If $A,B,C$ are three \emph{consecutive} points among the six, a chord separates them from $D,E,F$, so the two triangles (which are convex) lie on opposite sides of it and are disjoint. There are $6$ such splits (the blocks $123,234,\dots,612$).
\item Otherwise, going round the circle, the labels alternate at least twice, so there are two points of $DEF$ with points of $ABC$ on both arcs between them. The chord joining those two points of $DEF$ is a side of triangle $DEF$, and it separates two vertices of triangle $ABC$, so it crosses triangle $ABC$. The triangles intersect.
\end{itemize}
So the triangles are disjoint with probability $\dfrac{6}{20}=\dfrac{3}{10}$, and they intersect with probability
\[
1-\dfrac{3}{10}=\boxed{\dfrac{7}{10}}.
\]
(The number $2012$ plays no role.)
